%=====================================================================
\chapter{Qualitative Methods and Coding}\label{ch:qualitative}
%=====================================================================
\locator{3}{Part~3 --- Research Designs for Computing}

\begin{outcomes}
By the end of this chapter you will be able to:
\begin{tight}
  \item \textbf{Design} an interview guide that produces analysable data.
  \item \textbf{Choose} between grounded theory and thematic analysis on
    principled grounds.
  \item \textbf{Code} qualitative data through open, axial and selective
    stages, with memos.
  \item \textbf{Judge} when inter-coder reliability is appropriate and when it
    is a category error.
  \item \textbf{Defend} a claim of saturation with evidence rather than
    assertion.
\end{tight}
\end{outcomes}

\begin{prereq}
\chref{philosophy} --- the interpretivist and critical realist positions, and
why importing positivist criteria into this work is a mistake.
\chref{casestudy} shares much of the data collection machinery.
\end{prereq}

\begin{keyterms}
semi-structured interview $\bullet$ interview guide $\bullet$ probe $\bullet$
focus group $\bullet$ purposive sampling $\bullet$ theoretical sampling
$\bullet$ saturation $\bullet$ grounded theory $\bullet$ constant comparison
$\bullet$ open, axial and selective coding $\bullet$ memo $\bullet$ thematic
analysis $\bullet$ codebook $\bullet$ inter-coder reliability $\bullet$
Krippendorff's alpha $\bullet$ member checking $\bullet$ reflexivity
\end{keyterms}

\section{What Qualitative Work Is For}

Quantitative methods tell you how much and how often. Qualitative methods tell
you what something \emph{is}, how it works, and why people do what they do.
When the categories themselves are unknown --- when you cannot yet write the
questionnaire because you do not know what to ask --- there is no quantitative
substitute.

In computing this arises constantly. Why do teams abandon a practice they
believe in? What does ``technical debt'' actually mean to the people who use
the phrase? How do reviewers decide what deserves a comment? These are not
questions with numeric answers.

\begin{alertbox}[The most common error is importing the wrong criteria]
A quantitatively trained reader asks a qualitative study: is the sample
representative? What is the margin of error? Have you controlled for
confounders?

\medskip
These are category errors. A study seeking to understand meaning in context is
not estimating a population parameter, and 14 purposively chosen participants
are not a failed attempt at 400 random ones. The right questions are different
--- and \reg{PhD Regs 2024}{\S14.3(A)} asks them: does the sample produce the
\emph{type} of knowledge needed to understand the structures involved, is the
description detailed enough to interpret, are different sources compared and
contrasted, and does the analysis move from quotation to interpretation of
significance?

\medskip
Learn to ask those, and to answer them when asked the others.
\end{alertbox}

\section{Interviews}

\begin{center}
\small
\begin{tabular}{@{}L{2.8cm}L{4.6cm}L{6.9cm}@{}}
\toprule
\textbf{Type} & \textbf{Structure} & \textbf{Use when} \\
\midrule
Structured    & Fixed questions, fixed order
              & You know the categories and want comparability. Close to an
                orally administered survey \\
Semi-structured & A guide of topics, freely ordered, with probing
              & The default in computing research. Comparable across
                participants, and still able to follow the unexpected \\
Unstructured  & A single opening prompt and listening
              & Genuinely exploratory work, or life-history style research \\
\bottomrule
\end{tabular}
\end{center}

\begin{templatebox}[C9 --- Interview guide and codebook]
\textbf{Interview guide}

\medskip
\begin{tabular}{@{}L{3.6cm}L{9.9cm}@{}}
Opening (rapport, consent, recording permission) & \fillline{9.6cm} \\[7pt]
Warm-up: role and context & \fillline{9.6cm} \\[7pt]
Topic 1 --- addresses RQ\underline{\hspace{5mm}} & \fillline{9.6cm} \\[7pt]
\quad probes & \fillline{9.6cm} \\[7pt]
Topic 2 --- addresses RQ\underline{\hspace{5mm}} & \fillline{9.6cm} \\[7pt]
\quad probes & \fillline{9.6cm} \\[7pt]
Critical incident prompt & \fillline{9.6cm} \\[7pt]
Closing: anything I did not ask? & \fillline{9.6cm} \\[7pt]
\end{tabular}

\medskip
\textbf{Codebook} --- one row per code, maintained as coding proceeds

\medskip
\begin{tabular}{@{}L{2.6cm}L{3.4cm}L{3.4cm}L{3.6cm}@{}}
\textbf{Code} & \textbf{Definition} & \textbf{Include when} & \textbf{Exclude when} \\
\fillline{2.4cm} & \fillline{3.2cm} & \fillline{3.2cm} & \fillline{3.4cm} \\[7pt]
\fillline{2.4cm} & \fillline{3.2cm} & \fillline{3.2cm} & \fillline{3.4cm} \\[7pt]
\fillline{2.4cm} & \fillline{3.2cm} & \fillline{3.2cm} & \fillline{3.4cm} \\[7pt]
\end{tabular}
\end{templatebox}

\begin{conceptbox}[Four habits that separate usable interviews from unusable ones]
\begin{tightnum}
  \item \textbf{Ask about specific incidents, not general policy.} ``Walk me
    through the last code review you did'' produces data. ``What is your code
    review process?'' produces the official answer, which is often not what
    happens.
  \item \textbf{Tolerate silence.} The three seconds after an answer is where
    the qualification comes. Interviewers who fill it lose their best material.
  \item \textbf{Probe with the participant's own words.} ``You said it was
    messy --- what did messy look like that day?'' Never substitute your term
    for theirs; that is how the researcher's categories replace the
    participant's.
  \item \textbf{Record and transcribe.} Notes reconstruct what you already
    believed. Budget four to six hours of transcription per interview hour, or
    a transcription tool plus correction.
\end{tightnum}
\end{conceptbox}

\section{Sampling and Saturation}

Qualitative sampling is \term{purposive}: participants are chosen because they
can inform the question. In grounded theory it becomes \term{theoretical
sampling} --- each round of participants is chosen based on what the previous
round revealed, which means you cannot fix the sample in advance and should not
pretend to.

\begin{definitionbox}[Saturation]
\term{Saturation} is the point at which additional data stops producing new
codes, categories or theoretical insight.

\medskip
It is the standard justification for a qualitative sample size, and it is
routinely asserted rather than demonstrated~\cite{saunders2018saturation}. It is
also frequently misapplied: saturation belongs to grounded theory's iterative
logic, where sampling and analysis alternate. A study that conducts all 15
interviews before analysing any of them cannot know where saturation occurred.
\end{definitionbox}

\begin{center}
\small
\begin{tabular}{@{}L{4.4cm}L{9.9cm}@{}}
\toprule
\textbf{Instead of asserting it} & \textbf{Show it} \\
\midrule
``Saturation was reached at 14''
  & Report new codes per interview: a table or curve showing 9 new codes in
    interview 1, 2 in interview 10, 0 in interviews 13 and 14 \\
``No new themes emerged''
  & State the stopping rule \emph{in advance}: two consecutive interviews
    yielding no new codes \\
Silence on the point
  & If sampling was constrained by access rather than saturation, say so. It is
    a limitation, not a disqualification \\
\bottomrule
\end{tabular}
\end{center}

\section{Two Analytic Approaches}

\subsection{Grounded theory}

Grounded theory builds a theory from data, through iterative coding and
constant comparison, with sampling directed by the emerging theory. It is
demanding, and it is invoked far more often than it is done.

\begin{paperbox}[Grounded theory as it is actually practised in our field]
Klaas-Jan Stol, Paul Ralph and Brian Fitzgerald, \emph{Grounded theory in
software engineering research: a critical review and guidelines}, ICSE
2016~\cite{stol2016grounded}.

\medskip
\textbf{Why it matters.} They reviewed papers claiming grounded theory and
found most used it as a synonym for ``we coded some interviews'' --- without
theoretical sampling, without constant comparison, and without producing a
theory. If you are going to claim grounded theory, read this first and make
sure you can meet the claim.

\medskip
\textbf{What to extract.} Which variant you are following --- Glaserian,
Straussian or Charmaz's constructivist~\cite{charmaz2014constructing} --- and
their reporting guidelines. The variants differ on questions that matter,
particularly whether you may review the literature first.
\end{paperbox}

\begin{center}
\small
\begin{tabular}{@{}L{2.6cm}L{5.4cm}L{6.3cm}@{}}
\toprule
\textbf{Stage} & \textbf{What you do} & \textbf{Output} \\
\midrule
Open coding
  & Line-by-line, naming what is happening. Codes stay close to the data ---
    gerunds help: ``deferring the refactor'', not ``technical debt''
  & Many codes, deliberately unrefined \\
Axial coding
  & Relate codes to each other: conditions, actions, consequences
  & Categories with structure \\
Selective coding
  & Identify the core category and relate everything to it
  & A theory with a centre \\
Memoing
  & Throughout: write your thinking as it happens, dated
  & The audit trail, and most of the eventual discussion chapter \\
\bottomrule
\end{tabular}
\end{center}

\begin{alertbox}[Memos are not optional]
Memoing is what distinguishes grounded theory from coding. A memo records why
you created a code, how two categories seem to relate, what surprised you, and
what you now want to sample for.

\medskip
Without memos, the theory appears at the end with no traceable derivation, and
neither you nor an examiner can reconstruct how you got there. With them, the
audit trail is complete and the discussion chapter is half-written. Date every
memo; they are also your evidence of when an idea occurred, which matters if
anyone asks whether it preceded or followed the data.
\end{alertbox}

\subsection{Thematic analysis}

Braun and Clarke's six phases~\cite{braun2006using} are more accessible and, for
most graduate projects with a defined research question, more appropriate than
grounded theory. Thematic analysis identifies patterns; it does not claim to
build a theory, and it does not require theoretical sampling.

\begin{center}
\small
\begin{tabular}{@{}L{0.6cm}L{3.4cm}L{9.8cm}@{}}
\toprule
1 & Familiarisation & Read everything, repeatedly, before coding anything \\
2 & Initial codes   & Systematic coding across the whole dataset \\
3 & Searching for themes & Collate codes into candidate themes \\
4 & Reviewing themes & Check themes against coded extracts and the whole
                       dataset; some collapse, some split \\
5 & Defining themes & Name each and write what it is and is not. If you cannot
                     define it in two sentences, it is not a theme \\
6 & Writing up     & Analytic narrative with data extracts as evidence, not as
                     decoration \\
\bottomrule
\end{tabular}
\end{center}

\begin{pitfall}[Themes do not ``emerge'']
The word is a habit and it hides a claim. Themes are constructed by an analyst
making choices --- which codes to group, what to call them, what to leave out.
Saying they emerged implies they were sitting in the data waiting to be found,
which removes the analyst from a process they in fact drove.

\medskip
Write instead: ``we constructed four themes''. Then say who did the coding,
what their background was, and how disagreements were resolved. That is
reflexivity, and it is what \S14.3(A) means by treating subjective perceptions
as knowledge in their own right --- including the researcher's.
\end{pitfall}

\section{Reliability, and When It Does Not Apply}

Whether two coders should agree is a genuinely contested question, and the
answer depends on your paradigm (\chref{philosophy}).

\begin{center}
\small
\begin{tabular}{@{}L{4.0cm}L{10.3cm}@{}}
\toprule
\textbf{Position} & \textbf{Argument} \\
\midrule
Reliability is required
  & If codes are meant to be descriptive categories applied consistently, two
    trained coders should apply them alike. Report Cohen's $\kappa$ (two
    coders) or Krippendorff's $\alpha$ (any number, any
    scale)~\cite{krippendorff2018content} \\
Reliability is a category error
  & Under a constructivist position, interpretation is the researcher's
    contribution. Demanding that two people converge asks for a consistency the
    method does not claim, and rewards shallow, safe codes \\
\bottomrule
\end{tabular}
\end{center}

\begin{conceptbox}[A defensible position for a thesis]
Split the difference along the analysis, because the two halves have different
characters.

\medskip
\textbf{Descriptive coding} --- ``mentions a code review tool'', ``describes a
deadline conflict'' --- should be reliable, and reporting $\alpha$ is
appropriate and cheap.

\medskip
\textbf{Interpretive theme construction} --- what the pattern means --- is
argued rather than measured. Here the appropriate evidence is a documented
audit trail, negative case analysis (actively seeking data that contradicts
your theme), peer debriefing, and member checking with participants.

\medskip
Say in the methods chapter which parts of the analysis are which, and apply the
right standard to each. That is a stronger position than either extreme, and it
survives examination by a committee containing both kinds of reader.
\end{conceptbox}

\begin{traditions}{qualitative inquiry}
\tradEMP{Regards it as hypothesis-generating: the interviews tell you what to
measure, and the measurement comes later. Useful, and an incomplete account of
what qualitative work contributes.}
\tradDSR{Supplies the problem characterisation before design and the
naturalistic formative evaluation after --- what did people do with the
artefact that you did not anticipate?}
\tradQUAL{Its home. Meaning is constructed, context is constitutive, the
researcher is part of the instrument, and rigour comes from transparency and
audit rather than from replication.}
\tradFORM{Rarely engaged, but qualitative study of how practitioners reason
about formal methods is one of the more revealing literatures about why formal
techniques are not adopted.}
\end{traditions}

\begin{lab}[Code one interview three ways]
Obtain a transcript --- conduct one interview yourself, or use a published
transcript from an open dataset.

\begin{tightnum}
  \item Open-code it line by line using gerunds. Aim for 30 or more codes.
  \item Have a colleague code the same transcript independently with your
    codebook. Compute agreement on the descriptive codes only.
  \item Discuss every disagreement. Revise the codebook definitions --- the
    include and exclude columns of template C9 --- rather than negotiating
    individual segments.
  \item Group your codes into three or four candidate themes, and write a
    two-sentence definition of each saying what it is \emph{and is not}.
  \item Write one memo, dated, on the most surprising thing in the transcript.
\end{tightnum}

\medskip
Use Taguette (open source) or QualCoder if you want tool support; a spreadsheet
is entirely adequate for one transcript.
\end{lab}

\begin{checkpoint}
\begin{tightnum}
  \item Why does asking about a specific recent incident produce better data
    than asking about general practice?
  \item A thesis reports ``saturation was reached after 12 interviews''. What
    evidence should accompany that claim?
  \item Under what circumstances is demanding inter-coder reliability a
    category error, and under what circumstances is omitting it a weakness?
  \item Rewrite: ``Three themes emerged from the data.''
\end{tightnum}
\end{checkpoint}

\begin{chaptersummary}
\begin{tight}
  \item Qualitative methods answer what, how and why when the categories are
    not yet known. Judging them by sampling representativeness is a category
    error --- \S14.3(A) supplies the right criteria.
  \item Semi-structured interviews are the field's default. Ask about specific
    incidents, tolerate silence, probe in the participant's words, record and
    transcribe.
  \item Sample purposively; in grounded theory, theoretically. Demonstrate
    saturation with new-codes-per-interview, or admit access constrained you.
  \item Grounded theory requires theoretical sampling, constant comparison,
    memoing and an actual theory. Do not claim it otherwise.
  \item Thematic analysis is usually the right choice for a graduate project
    with a defined question.
  \item Themes are constructed, not emergent. Say who constructed them.
  \item Apply reliability to descriptive coding; apply audit trail, negative
    case analysis and member checking to interpretation. State which is which.
\end{tight}
\end{chaptersummary}

\begin{reviewq}
  \item Stol and colleagues found most software engineering papers claiming
    grounded theory did not practise it. What incentives produce this, and what
    would a reviewer have to check to detect it?
  \item Negative case analysis asks you to actively seek data contradicting
    your theme. Compare this with falsification in \chref{philosophy}: are they
    the same epistemic move in different vocabularies, or genuinely different?
  \item Automated coding using large language models is increasingly proposed
    for qualitative analysis. Given \reg{PhD Regs 2024}{\S32(IV)}, which
    prohibits AI from ``analyzing and interpreting data'', what --- if anything
    --- may legitimately be automated here? (\chref{genai} takes this further.)
  \item Member checking returns findings to participants for comment. Give a
    case where a participant's disagreement should change your analysis and one
    where it should not, and state the principle that distinguishes them.
\end{reviewq}
